%% kapa: exam=1 round=22 year=2011 period=1 form=A subject=civil_law subject_ko=민법 seq=40 number=40
%% source: Q-Net 감정평가사 기출문제 (cst00309) artlSeq=5036790 file=extracted/1008/5036790/2011년 제22회 감정평가사 1차시험 기출문제/1교시.hwp
%% conversion: kordoc 4.24.51 -> tools/gen1.py (auto); review=proofread-sonnet; flags=none
\begin{question}{40}
동산질권에 관한 설명으로 옳지 않은 것은?
\begin{choices}
  \item 동산질권은 점유개정으로도 성립할 수 있다.
  \item 양도할 수 없는 동산은 질권의 목적물로 될 수 없다.
  \item 질권자는 질물로부터 생기는 천연과실을 수취하여 다른 채권보다 먼저 그의 채권의 변제에 충당할 수 있다.
  \item 질권자는 피담보채권 전부를 변제받을 때까지 질물을 유치할 수 있다.
  \item 질권자가 질물을 유치하고 있더라도 피담보채권의 시효는 진행한다.
\end{choices}
\end{question}
\answer{1}
